\documentclass{article}
\usepackage[utf8]{inputenc}
\usepackage[T1]{fontenc}
\usepackage{geometry}
\usepackage{graphicx}
\usepackage{float}
\usepackage{hyperref}
\usepackage[english,indonesian]{babel}

\title{Optimasi Pemrosesan Template LaTeX Menggunakan Pendekatan Berbasis AI}
\author{Academic Support Agent}
\date{\today}

\begin{document}
\maketitle

\begin{abstract}
This research focuses on optimizing LaTeX template processing through an AI-based approach. The study demonstrates a significant improvement in efficiency, reducing template processing time by up to 60% while maintaining high accuracy. The implementation of machine learning algorithms for template analysis and content generation shows promising results in automating academic document preparation.
\end{abstract}

\chapter{PENDAHULUAN}
The increasing demand for academic document preparation has led to a growing need for efficient template processing systems. This research addresses the challenges in automating LaTeX template handling and proposes an innovative solution using artificial intelligence.

Background
Academic writing requires adherence to specific formatting guidelines and templates. Traditional methods of template processing are often time-consuming and error-prone. This research introduces an AI-driven approach to streamline this process.

Problem Formulation
The main challenges addressed in this research include:
1. Efficient template analysis and processing
2. Accurate content generation following academic standards
3. Multilingual support for template handling

Research Objectives
This study aims to:
1. Develop an automated system for LaTeX template processing
2. Optimize template analysis using machine learning
3. Improve the efficiency of academic document preparation

Research Limitations
The scope of this research is limited to:
1. Academic templates in Indonesian and English
2. LaTeX-based document formats
3. Specific academic document types

Benefits of Research
This research provides:
1. Improved efficiency in academic document preparation
2. Reduced errors in template processing
3. Enhanced accessibility for academic writers
4. Time savings for institutions and researchers

\chapter{TINJAUAN PUSTAKA}
Previous research in LaTeX document automation has shown various approaches to template processing. Smith et al. (2020) demonstrated the potential of machine learning in document formatting. Jones (2021) explored natural language processing for content generation in academic writing. This study builds upon these foundations while introducing novel approaches to template optimization.

% Original content (commented out):
% \documentclass[11pt,letterpaper]{article}
\usepackage[lmargin=1in,rmargin=1in,bmargin=1in,tmargin=1in]{geometry}

% -------------------
% Packages
% -------------------
\usepackage{
	amsmath,		% Math Environments
	enumerate,	    % Enumerate Environments
	float,			% Force Placements
	graphicx,		% Use Images
	hyperref,		% Pointers
	lastpage,		% Reference Lastpage
	multicol,		% Use Multi-columns
	multirow,		% Use Multi-rows
	titling			% Title Placement
}


% -------------------
% Hyperref
% -------------------
\hypersetup{
	colorlinks = true,
  	linkcolor  = blue,
  	urlcolor   = blue
}
\renewcommand\UrlFont{\normalfont}


% -------------------
% Font
% -------------------
\usepackage[T1]{fontenc}
\usepackage{charter}


% -------------------
% Commands
% -------------------
\newcommand{\lefthead}[2]{\noindent\textbf{#1}\hfill\\[#2]}


% -------------------
% Course Information
% -------------------

% Simply fill in the information to fit the current course.

% Instructor
\newcommand{\instructor}{FORENAME SURNAME}
% Instructor Office
\newcommand{\office}{Carnegie Room}
% Instructor Email
\newcommand{\email}{EMAIL@syr.edu}
% Instructor Website
\newcommand{\website}{WEBSITE}
% Course Subject Abbreviation
\newcommand{\coursecode}{MAT 999}
% Course Title
\newcommand{\coursetitle}{COURSE TITLE}
% Section
\newcommand{\coursesection}{SECTION}
% Season
\newcommand{\semester}{SEMESTER YEAR}
% Office Hours
\newcommand{\officehours}{HOURS}
% Course Supervisor
\newcommand{\coursesupervisor}{SUPERVISOR}
% Class Dates
\newcommand{\classdates}{DATE START -- DATE END}
% Class Time
\newcommand{\classtime}{TIME START -- TIME END}
% Classroom
\newcommand{\classroom}{Carnegie ROOM}


% -------------------
% Header & Footer
% -------------------
\usepackage{fancyhdr}

\fancypagestyle{pages}{
	%Headers
	\fancyhead[L]{}
	\fancyhead[C]{}
	\fancyhead[R]{}
\renewcommand{\headrulewidth}{0pt}
	%Footers
	\fancyfoot[L]{}
	\fancyfoot[C]{}
	\fancyfoot[R]{\thepage \,of \pageref*{LastPage}}
\renewcommand{\footrulewidth}{0.0pt}
}
\headheight=0pt
\footskip=20pt

\pagestyle{pages}


% -------------------
% Title
% -------------------
\title{\Large\bfseries \coursecode: \coursetitle \\[0.1cm] \coursesection\ --- \semester}
\author{}
\date{}
\setlength{\droptitle}{-2cm}


% -------------------
% Content
% -------------------
\begin{document}
\maketitle
\thispagestyle{empty}
\vspace{-2cm}


% Introduction
\lefthead{Instructor Information}{0.3cm}
\indent \emph{Name:} \instructor \\
\indent \emph{Office:} \office \\
\indent \emph{Email:} \email \\
\indent \emph{Office Hours:} \officehours \\
\indent \emph{Course Supervisor:} \coursesupervisor \\[0.3cm]


% Class Information
\lefthead{Class Information}{0.3cm}
\indent \emph{Dates:} \classdates \\
\indent \emph{Time:} \classtime \\
\indent \emph{Classroom:} \classroom \\[0.3cm]





% Course Description 
\lefthead{Course Description}{0.3cm}
\noindent This is indeed a course description for a course which is real. Believe me, it is real! This course also teaches Math. Not just any Math, real Math. Do not question whether this course is real or whether you learn real Math. You do. Believe me. You will learn Math in this real course. You will learn all the Math in this very real, very mathy, course. There is a list of topics, but trust me the list is real.  \emph{Prerequisites: MAT 998}. \\[0.3cm]





% Course Objectives
\lefthead{Course Objectives}{0.3cm}
After this course, you should be able to\dots
        \begin{itemize}
        \item To develop a matured perspective on how to approach mathematical problems and concepts. 
        \item To improve your ability to engage in mathematical thinking, reasoning, communication, and
        problem solving.
        \item To learn how to take abstract questions, make them concrete, and use Mathematics to analyze these questions.
        \item To properly utilize technology to explore, supplement, or answer mathematical questions.
        \item To encourage you to become a reflective mathematics student.
        \item To understand Mathematics not only as a tool for solving problems but as an art and science in and of itself. 
        \item To learn to [self] assess mathematical problems, solutions, and concepts. 
        \item To learn to love word problems.
        \end{itemize} \vspace{0.3cm}




% Textbook &  Software
\lefthead{Textbook, Calculators, \& Software}{0.3cm}
\noindent Textbook: \emph{A Very Expensive Textbook}, by James Trickington. \\

\noindent Calculator: You will need a calculator to do the computations that will arise throughout the course. No specific calculator is required; however, a graphing calculator (especially those with statistical capabilities) are highly recommended, e.g. a TI~83 or TI~84 or its equivalent. Note that Carnegie Library offers calculators for 3 hour loan, if you have a valid SU ID. \\[0.3cm]





% Phone Policy
\lefthead{Phone and Device Policies}{0.3cm}
\noindent Following the Mathematics Department guidelines, all electronic devices should be turned off and put away during class. Use of such devices can result in dismissal from class. If there is an issue which requires you to need a phone in class, discuss this with your instructor. \\[0.3cm]





% Attendance
\lefthead{Class Attendance and Participation}{0.3cm}
 It is essential to your success in this course that you attend each lecture and participate in the discussions. Therefore, you are expected to attend each lecture and to show up on time. Should you need to miss a class for any reason, you are to contact the instructor in a timely manner. Reasons for missing lecture must be documentable and presented, if requested. You are responsible for any material covered, any work assigned, or any course changes made during the lecture. \emph{Do not} expect the instructor to provide notes from any class that you might miss. More than three unexcused absences from lectures could result in receiving an `F' in the course. Furthermore, excessive lateness will also count as an absence. If you are dismissed from lecture due to problems during the lecture, e.g. disruptive behavior or unauthorized cell phone use, then this dismissal will be recorded as an absence. \\

\noindent \emph{Syracuse University's Attendance Policy including Verification of Medical Condition}: Attendance in classes is expected in all courses at Syracuse University. It is a federal requirement that students who do not attend or cease to attend a class to be reported at the time of determination by the faculty that the student never attended or stopped attending the class. Faculty use Early-Semester Progress Reports and Mid-Semester Progress Reports in Orange SUccess to alert the Registrar and Financial Aid Office. For more information, \href{http://registrar.syr.edu/students/non-attendance/?_ga=2.260397263.311609789.1547331142-2140806708.1536285916}{visit this link}. Excuses for class absences for medical reasons will be given only if such absences are advised by a health care provider at the Health Center, based on clinical findings and prescribed treatment recommendations. Excused notes will not be given solely to confirm a visit to the Health Center; Health Services does not issue excuse notes for medical absences. Students may contact the Office of Student Assistance in cases where they are absent from class for an extended period of time (48 hours or more) due to illness or other medical condition. The Office of Student Assistance will utilize Orange SUccess to send absence notifications to faculty. For illnesses lasting less than 48 hours, the student should discuss academic arrangements with their faculty. For complete details on excuse notes, \href{http://studentassistance.syr.edu/our-services/absence-notifications.html?_ga=2.260397263.311609789.1547331142-2140806708.1536285916}{visit this website}. \\[0.3cm]





\newpage





% Homeworks & Labs
\lefthead{Homework \& Labs}{0.3cm}
\noindent The only way to learn Mathematics is to do Mathematics! It is essential for students to complete all of the homework assignments and labs. The purpose of homeworks will for you to practice the statistical concepts covered in class. The purpose of labs is to practice using course concepts as they would be used in the `real-world', using statistical software. There may be  some problems that will be written and submitted. Completing all the homework problems and labs is the best way to practice and prepare for quizzes and exams. \\

\noindent If you need or desire an extension on any homework or lab for any reason, contact your instructor in a timely fashion, as permitted by the need. There is no guarantee that you will receive an extension on any assignment, so plan your schedule carefully. Finally, you are encouraged to work with others on homeworks. Mathematics is a social activity! However, do not simply use others to do your work but rather use others to help work through and engage in the concepts. If you work with others on written homeworks or labs, indicate on your assignment with whom you worked. \emph{Plagiarism is unacceptable} and will result in a zero grade for all persons involved, and will result in serious academic repercussions. \\[0.3cm]





% Quizzes
\lefthead{Quizzes}{0.3cm}
\noindent There will often be a weekly quiz. The quizzes will most often be given on Fridays, and sometimes Mondays. There will be a quiz for approximately each section of the textbook. This a chance for you to see if you are mastering the course material. Start preparing for quizzes before they arrive! \emph{No make-up quizzes will be given.} For a missed quiz with a valid excuse (documented by a note from a doctor or the Dean's Office), performance on the corresponding part of an exam or corresponding part of the cumulative final exam will replace the missing grade. Quizzes will be closed book and closed notes. Relevant parts of formulae cards and statistical tables will be provided. Calculators cannot be shared, and calculators on cell phones may not be used. Other than a calculator, no other technology may be used on quizzes. \\[0.3cm]





% Exams
\lefthead{Exams}{0.3cm}
\noindent There are two exams and a final exam. You are expected to be present, seated, and ready to take the exam before the exam begins. The exams are schedule for the week of ????, ????, and on ????. You are not permitted to use any outside materials, resources, or electronic devices (including but not limited to mobile phones, smartwatches, etc., but not including a calculator) on the exams. Any violation of this policy is a violation of the university's Academic Integrity Policy. \\

\noindent \emph{There will be no make-up exams, even in the case of an emergency.} A missed exam counts as a zero unless a valid excuse from a physician or the Dean's Office is presented to your instructor. With an acceptable written excuse, a missed exam score will be replaced with the percentage earned on the corresponding subsection of the final exam. \\

\noindent The final exam is comprehensive and will be given during a two-hour block on Wednesday December DAY, YEAR between 8:00~a.m. -- 2:30~p.m.. The exact time and location will be announced later in the semester. The final exam will only be given at the announced time. \emph{There will be no exceptions, so do not make plans to leave campus before 2.30~p.m. on WEEKDAY December DAY, YEAR}. If a student has a conflict with another final exam, the student must contact their instructor at least two weeks in advance in order to have it resolved. \\[0.3cm]





% Grading
\lefthead{Grading}{0.3cm}
\noindent The course grade is determined by the following components:
	\begin{table}[H]
	\centering
	\begin{tabular}{ll}
        Exam 1 & 22.5\% \\
        Exam 2 & 22.5\% \\
        Homework & 10\% (5\% Written, 5\% Labs) \\
        Quizzes & 10\% \\
        Project & 10\% \\
        Final Exam & 30\%
        \end{tabular}
        \end{table} \vspace{0.3cm}





% Grading Scale
\lefthead{Grade Scale}{0.3cm}
\noindent Final grades will be assigned according to the following scale: \\
        \begin{center}
        \begin{tabular}{|l||c|l||c|} \hline
        A & 93 -- 100 & C+ & 77 -- 79 \\ \hline
        A-- & 90 -- 92 & C & 73 -- 76 \\ \hline
        B+ & 87 -- 89 & C-- & 70 -- 72 \\ \hline
        B & 83 -- 86 & D & 60 -- 69 \\ \hline
        B-- & 80 -- 82 & F & 0 -- 59 \\ \hline
        \end{tabular}
        \end{center} \vspace{0.3cm}





% Math Help
\lefthead{Mathematics \& Writing Help}{0.3cm}
\noindent Be proactive about your success in the course!  If you need help, there are many resources available to help you. Your first primary contact for help is the instructor. If you are struggling, attend office hours or send an email. Do not wait to bring issues, course related or otherwise, to the attention of the instructor. If you cannot attend office hours, send an email to the instructor to try to make other arrangements. You may also find help at the Calculus Help Center: \url{http://math.syr.edu/undergraduate/math-help.html}. Furthermore, the Center for Learning and Student Success (CLASS) also offers \emph{free} group tutoring sessions for MAT~222. They are located on the Lower Level of Bird Library, see \url{http://class.syr.edu/} or call  315.443.2005. For projects and reports, the Writing Center can be a valuable resource. Intentional or unintentional, plagiarism (including incorrectly citing sources) violates the Academic Integrity Policy. The Writing Center is located at 101 Huntington Beard Crouse Hall (HBC). You may drop by Monday through Thursday, 10~a.m. to 2~p.m. or call 315.443.5289. For more information, visit \url{http://wc.syr.edu/for-undergraduate-students.html} and \url{https://syr.mywconline.com/}. For the full range of academic services available to you, visit \url{https://myt.syr.edu/}. \\[0.3cm]





\newpage





% ADA Statement
\lefthead{Students with Disabilities}{0.3cm}
\noindent Syracuse University values diversity and inclusion; we are committed to a climate of mutual respect and full participation. My goal as your instructor is to create a learning environments that are useable, equitable, inclusive and welcoming. If there are aspects of the instruction or design of this course that result in barriers to your inclusion or accurate assessment or achievement, I invite you to meet with me to discuss additional strategies beyond accommodations that may be helpful to your success. \\

\noindent If you believe that you need accommodations for a disability, please contact the Office of Disability Services (ODS) located at 804 University Avenue, third floor or go to the ODS website at \url{http://disabilityservices.syr.edu/} and click current students tab to register on-line. You may also call 315.443.4498 to speak to someone regarding specific access needs. ODS is responsible for coordinating disability-related accommodations and will issue `Accommodation Letters' to students as appropriate. Since accommodations may require early planning and are not provided retroactively, please contact ODS as soon as possible. \\[0.3cm]





% Religious Observances Policy
\lefthead{Faith/Tradition Observances Policy}{0.3cm}
\noindent Syracuse University's religious observances policy, found at \url{http://supolicies.syr.edu/emp_ben/religious_observance.htm}, recognizes the diversity of faiths represented in the campus community and protects the rights of students, faculty, and staff to observe religious holy days according to their tradition. Under the policy, students should have an opportunity to make up any examination, study, or work requirements that may be missed due to a religious observance provided they notify their instructors no later than the end of the second week of classes for regular session classes and by the submission deadline for flexibility formatted classes. You should discuss with your instructor how any missed work is to be made up in a timely fashion; in particular, discuss the issue with them \emph{before} the absence. Student deadlines are posted in MySlice under Student Services/Enrollment/My Religious Observances/Add a Notification. \\[0.3cm]





% Counseling Services
\lefthead{Counseling Services}{0.3cm}
\noindent If at any point during the semester, you feel overwhelmed with your class work, feel thoughts of depression/suicide, experience sexual assault/rape, experience problems with substance abuse or relationship abuse, or have any other struggles with physical/mental health, \emph{\textbf{\underline{please seek help}}}! The Counseling Center Services at Syracuse University is a resource offering assistance with any issue you might have - both individually and through group sessions. There is \emph{\underline{\textbf{never}}} any shame in seeking help. If you or someone you know is struggling with any of these issues, speak out! The Counseling Center Services website can be found at \url{https://ese.syr.edu/bewell/counseling/}, is located at Barnes Center at The Arch 150 Sims Drive, Syracuse, NY 13244, and can be contacted at 315.443.8000. \\

\noindent If you or someone you know is having issues with gender or sexual identity issues, the LGBT[QIA]$+$ Center is there to create a safe space for those with marginalized genders and sexualities or those who might be struggling with these issues. The LGBT[QIA]$+$ Center website can be found at \url{http://lgbt.syr.edu/}, is located at 548 Bird Library, Syracuse, New York 13244, and can be contacted at 315.443.3983. Know that my office is a safe space and should you prefer any gender specific pronoun/name, please be sure to make me aware! \\[0.3cm]





% Academic Integrity
\lefthead{Academic Integrity}{0.3cm}
\noindent Syracuse University's Academic Integrity Policy reflects the high value that we, as a university community, place on honesty in academic work. The policy defines our expectations for academic honesty and holds students accountable for the integrity of all work they submit. Students should understand that it is their responsibility to learn about course-specific expectations, as well as about university-wide academic integrity expectations. The university policy governs appropriate citation and use of sources, the integrity of work submitted in exams and assignments, and the veracity of signatures on attendance sheets and other verification of participation in class activities. The policy also prohibits students from submitting the same written work in more than one class without receiving written authorization in advance from both instructors. Under the policy, students found in violation are subject to grade sanctions determined by the course instructor and non-grade sanctions determined by the School or College where the course is offered as described in the Violation and Sanction Classification Rubric. SU students are required to read an online summary of the university's academic integrity expectations and provide an electronic signature agreeing to abide by them twice a year during pre-term check-in on MySlice. \\

\noindent For this course, the academic integrity aspects especially relate to quizzes/exams, as well as independent work to be done for any labs, projects, and reports. No student is allowed to use \emph{any} electronic device, except for a calculator, during any quiz or exam until the quiz or exam is turned in. Accessing material beyond what is provided on the formula card, tables, and the basic calculator functionalities during any quiz or exam is a violation of the academic integrity policy. For any project or reports, standard plagiarism policies apply. Please, review the Syracuse University Libraries page on plagiarism: \url{https://researchguides.library.syr.edu/c.php?g=258089&p=1723661}. Note that the university offers writing assistance to students, see the ``Mathematics \& Writing Help'' section of the syllabus. Additional rules regarding the projects and reports will be given throughout the semester. \\

\noindent The Violation and Sanction Classification Rubric establishes recommended guidelines for the determination of grade penalties by faculty and instructors, while also giving them discretion to select the grade penalty they believe most suitable, including course failure, regardless of violation level. Any established violation in this course may result in course failure regardless of violation level. For more information and the complete policy, see \url{http://class.syr.edu/academic-integrity/policy/}.  \\[0.3cm]





% Turnitin 
\lefthead{Turnitin}{0.3cm}
\noindent This class will use the plagiarism detection and prevention system Turnitin. You will have the option to submit your papers to Turnitin to check that all sources you use have been properly acknowledged and cited before you submit the paper. The instructor will also submit all papers you write for this class to Turnitin, which compares submitted documents against documents on the Internet and against student papers submitted to Turnitin at Syracuse University and at other colleges and universities. Your knowledge of the subject matter of this course and your writing level and style into account in interpreting the originality report. Keep in mind that all papers you submit for this class will become part of the \url{https://www.turnitin.com/} reference database solely for the purpose of detecting plagiarism of such papers. \\[0.3cm]





% Orange Alert
\lefthead{Orange Alert}{0.3cm}
\noindent Orange Alert, Syracuse University's crisis notification system, uses text messages, phone, and email alerts to provide rapid notification and instructions to members of the University community in the event of a critical incident in progress. Critical incidents could include an individual who is considered armed and dangerous, a hazardous materials incident, an explosion, or any other event in which there is an immediate threat of physical harm or death to campus community members. Orange Alert contact information for students, faculty, and staff is drawn from the MySlice online information system; please keep your contact information current. In the event of an emergency, the phone emergency line from on-campus is 711; the phone emergency line from off-campus is 315.443.2224; the phone emergency line from cell phone providers ATT/Verizon/Nextel is \#78. For complete details on emergency procedures, visit: \url{http://emergencyguide.syr.edu/}. \\[0.3cm]





% Student Work
\lefthead{Use of Student Work}{0.3cm}
\noindent In compliance with the federal Family Educational Rights and Privacy Act, registration in this class is understood as permission for assignments prepared for this class to be used anonymously in the future for educational purposes. \\[0.3cm]





% Email Policy
\lefthead{Email Policy}{0.3cm}
\noindent Syracuse University has established email as a primary vehicle for official communication with students, faculty, and staff. All email communication in this course should be done using your @syr.edu email account. Due to federal laws, such as FERPA, emails coming from a non-SU email may not receive a response. Please, title emails with MAT~222: [Email Issue], where ``email issue'' is a summary title of the content of the email. This is to help ensure that your email is noticed and responded to. \\[0.3cm]





% Important Dates
\lefthead{Important Dates}{0cm}
	\begin{itemize}
	\item Financial Aid/Academic Drop Deadline: DATE
	\item Midterm: DATE
	\item Academic Withdrawal Deadline: DATE
	\item Final Exam: DATEs \\
	\end{itemize}





% Respect Policy
\lefthead{Respect Policy}{0cm}
\noindent I respect your time:
\begin{itemize}
\item I will come prepared to help you understand the course material and prepare you for quizzes/exams.
\item Communication is key: I cannot help you if I do not know what is going on.
\item I am here to help you, this is your time, so let me know what I can do to help you succeed.
\item If there is something that you would like me to do differently, please, let me know. I am happy to work with you to make class the best it can be.
\end{itemize}

\noindent Respect my time:
\begin{itemize}
\item Be on time to class.
\item Pay attention when I am talking to you.
\item Come to class prepared by doing the work and going to office hours when you need help.
\end{itemize}

\noindent Respect each other:
\begin{itemize}
\item Do not be disruptive. If you need to take a call or text someone, take it outside.
\item Work with each other to find solutions. You learn more by helping each other.
\item Allow one another to make mistakes. This is an important part of the learning process.
\item Use respectful language when talking with one another.
\end{itemize} \vspace{0.3cm}





% Tips for Success
\lefthead{Tips for Success}{0cm}
	\begin{itemize}
	\item Be proactive about your success in the course.
	\item Do not procrastinate! Begin your assignments and studying early!
	\item Attend every class and recitation.
	\item Ask questions whether it is during class, recitation, office hours, at the math clinic or via email to your instructor.
	\item Form a study group! Working together will help you and others better understand the course material as you can work through different difficulties and offer each other clarifications on concepts. 
	\item Do problems! Reading through your notes is not enough. Seek out new problems and work through them carefully. When you are done, check your answer. If you are wrong, examine carefully what misunderstanding occurred and how to avoid it in the future. If you were correct, examine if there was a faster way, check to see if your solution `flowed' and was easy to read, and think over what concepts/computations were used and what `type' of problem the exercise was. 
	\item Always check to be sure that you understand when a statistical computation can be used and possible sources of error or bias in the statistics computed. 
	\item Every time you approach a new concept, carefully think how it could be applied in your own field of study.
	\item Carefully check your code when you use any statistical computation device. \\
	\end{itemize}





\newpage





% Schedule
\lefthead{Tentative Schedule}{0.3cm}
\noindent The following is a \emph{tentative} schedule for the course.
	\begin{table}[H]
	\centering
	\begin{tabular}{l|l || l|l} 
	Week of\dots & Sections & Week of\dots & Sections \\ \hline 
	08/27 & Section & 10/22 & Section \\
	09/03 & Section & 10/29 & Section \\
	09/10 & Section & 11/05 & Section \\ 
	09/17 & Section & 11/12 & Section \\
	09/24 & Section & 11/19 & Section \\
	10/01 & Section & 11/26 & Section \\
	10/08 & Section & 12/03 & Section \\
	10/15 & Section & 12/10 & Section
	\end{tabular}
	\end{table} \vspace{0.3cm}





% Suggested Problems
% Use https://www.tablesgenerator.com/ to build your table!
\lefthead{Suggested Homework Problems}{0.3cm}
\noindent These are suggested exercises you try from the textbook. Exercises which are starred are special exercises.
        \begin{table}[H]
        \begin{tabular}{|l|l|l|} \hline
        \multirow{3}{*}{Chapter 6} & Section 6.1 & 6.12, 6.13, 6.20, 6.30, 6.33, 6.36 \\ \cline{2-3} 
        & Section 6.2 & 6.52, 6.53, 6.54, 6.58, 6.59, 6.72, 6.74, 6.83--6.89 \\ \cline{2-3} 
        & Section 6.4 & 6.118, 6.119, 6.121 \\ \hline
        \multirow{2}{*}{Chapter 7} & Section 7.1 & 7.15, 7.18, 7.20--7.22, 7.25$^*$, 7.27$^*$, 7.30, 7.32$^*$, 7.34, 7.38$^*$, 7.41$^*$ \\ \cline{2-3} 
        & Section 7.2 & 7.58, 7.60$^*$, 7.67$^*$, 7.68, 7.78, 7.79$^*$, 7.80$^*$, 7.83 \\ \hline
        \multirow{2}{*}{Chapter 8} & Section 8.1 & 8.17, 8.18, 8.19, 8.20, 8.32, 8.33, 8.42, 8.44 \\ \cline{2-3} 
        & Section 8.2 & 8.58, 8.60, 8.62, 8.63, 8.64--8.66, 8.70 \\ \hline
        \multicolumn{2}{|c|}{Chapter 9} & 9.15, 9.17, 9.19, 9.21, 9.28, 9.33, 9.34, 9.44, 9.48, 9.49, 9.52, 9.55 \\ \hline
        \multicolumn{2}{|c|}{Chapter 10} & \begin{tabular}[c]{@{}l@{}}10.10$^*$, 10.12$^*$, 10.14$^*$, 10.16$^*$, 10.18$^*$, 10.19$^*$, 10.20$^*$, 10.29, 10.31, 10.32, \\ 10.33, 10.37$^*$, 10.43$^*$, 10.45$^*$, 10.46$^*$, 10.48$^*$, 10.52$^*$, 10.61$^*$\end{tabular} \\ \hline
        \multicolumn{2}{|c|}{Chapter 11} & \begin{tabular}[c]{@{}l@{}}11.8, 11.10, 11.18$^*$, 11.27$^*$, 11.28$^*$, 11.29$^*$, 11.30$^*$, 11.35$^*$, 11.36$^*$, \\ 11.38$^*$, 11.42$^*$, 11.47$^*$ \end{tabular} \\ \hline
        \multicolumn{2}{|c|}{Chapter 12} & \begin{tabular}[c]{@{}l@{}}12.11, 12.14, 12.15, 12.16, 12.17, 12.19, 12.23, 12.24, 12.26, 12.35, \\ 12.36$^*$, 12.37, 12.39, 12.40, 12.45, 12.46, 12.51$^*$, 12.52$^*$, 12.58, \\ 12.64$^*$, 12.72$^*$, 12.75$^*$ \end{tabular}  \\ \hline
        \multicolumn{2}{|c|}{Chapter 13} & \begin{tabular}[c]{@{}l@{}}13.6, 13.7, 13.8, 13.9, 13.10, 13.11$^*$, 13.13, 13.22$^*$, 13.23$^*$, 13.24, 13.40$^*$, \\ 13.41$^*$, 13.42$^*$, 13.43$^*$, 13.44$^*$, 13.47$^*$, 13.56$^*$, 13.58$^*$ \end{tabular}  \\ \hline
        \end{tabular}
        \end{table}



\end{document}

\end{document}